\documentclass[a4paper]{ctexart}
\usepackage[margin=1in]{geometry}

\usepackage{caption} % 图表标题
\usepackage{float} % 浮动体
\usepackage{graphicx} % 插图片
\usepackage{subfig} % 子图片
\usepackage{longtable} % 长表格

\usepackage{amsmath} % 数学公式
\usepackage{amssymb} % 数学符号
\usepackage{amsthm} % 定理环境
\usepackage{bm} % 数学粗体
\usepackage{braket} % 左右矢

\usepackage{color} % 字体上色

\usepackage{fancyhdr} % 页眉页脚
\usepackage{lastpage} % 最后一页

\usepackage{ulem} % 标记线

\usepackage{cite} % 引用
%\usepackage[backend=biber,style=verbose-note]{biblatex}

\usepackage[colorlinks=true,linkcolor=RGB{200,0,0}]{hyperref}
%\addbibresource{refs.bib}
\definecolor{Nred}{RGB}{200,0,18}

\hypersetup{
	colorlinks=true,
	linkcolor=Nred,
}

\pagestyle{fancy} %页脚页眉
\headheight=13pt
\fancyhead[R]{$\mathfrak{Quantum\ Mechanics}$}
\fancyhead[L]{\kaishu 量子力学}
\fancyfoot[C]{\songti 第\thepage 页/共26页}%页脚设置
\renewcommand{\headrulewidth}{0.02cm}
\renewcommand{\footrulewidth}{0.02cm}
\fancypagestyle{first}{
    \fancyhead{}
    \fancyfoot{}
}

\ctexset{  %字体设置
    punct=quanjiao,
    autoindent=true,
    contentsname={目录},
    figurename={图},
    tablename={表},
    abstractname={摘\quad 要},
    bibname={参考文献},
    section={
        format=\large\heiti
    },
    subsection={
        format=\normalsize\heiti
    },
    subsubsection={
        format=\normalsize\kaishu
    }
}

\renewcommand{\i}{\mathrm{i}}
\renewcommand{\d}{\mathrm{d}}
\newcommand{\e}{\mathrm{e}}
\renewcommand{\P}{\mathbf{P}}
\newcommand{\T}{\mathbf{T}}
\newcommand{\C}{\mathbf{C}}

\numberwithin{equation}{section}

\newtheorem{example}{例}[section]
\newtheorem{theorem}{定理}[section]
\newtheorem{postulate}{假设}[section]
\newtheorem{lemma}{引理}[subsection]

\begin{document}
$\ $
\vspace{2cm}
\begin{center}
    \zihao{-1}\heiti 2026春量子力学（3学分）
\end{center}
\vspace{2cm}
\section*{课程信息}
\paragraph{上课时间地点}周二1-2节（仅单周），周四3-4节，二教411
\paragraph{习题课时间地点}周二1-2节（仅双周），二教410
\paragraph{授课教师}舒菁\quad jshu@pku.edu.cn
\paragraph{助教}商言菲、苟凯翔
\paragraph{课程参考教材}\quad \\
量子力学（刘玉鑫、曹庆宏）\qquad 量子力学现代教程（孙昌璞）\qquad 量子力学（曾谨言）\\
Introduction to QM (Griffins)\qquad Modern QM (J.J.Sakurai)\qquad Lectures on QM (S.Weinberg)
\newpage
\tableofcontents
\newpage
\section{量子力学的数学理论}
\subsection{“推导”薛定谔方程}
在20世纪初，人们发现，微观世界与经典力学相比，有三个不同之处。第一，能量是离散的（Planck）；第二，粒子具有波的某些性质，也就是波粒二象性（德布罗意）；第三，测量是有概率的（波恩）。

为了找到描述粒子运动的方程，一个自然的想法是从经典波动方程出发
$$\frac{\partial ^2u}{\partial t^2}-\frac{1}{c^2}\nabla^2 u=0\quad u=u(x,t)$$
这个方程有平面波解
$$u(x,t)=\psi(x,t)=e^{i\vec{k}\cdot\vec{x}-i\omega t}\quad \omega=c|\vec{k}|$$
这个平面波解又满足下述两个方程
$$i\frac{\partial}{\partial t}\psi=\omega\psi\quad -i\vec{\nabla}\psi=\vec{k}\psi$$
又根据德布罗意物质波的假定$E=\hbar\omega,\vec{p}=\hbar\vec{k}$，可以做下述的对应
$$\vec{p}\leftrightarrow -i\hbar\vec{\nabla}\quad E\leftrightarrow i\hbar\frac{\partial}{\partial t}$$
经典力学中的能量关系$E=\frac{p^2}{2m}+V(x)$便对应于薛定谔方程
\begin{equation}
    i\hbar\frac{\partial}{\partial t}\psi=-\frac{\hbar^2}{2m}\nabla^2\psi+V(x)\psi\quad\psi=\psi(x,t)
\end{equation}
$|\psi(x,t)|^2$表示$t$时刻，粒子出现在$x$位置的概率密度。换言之，在区域$Q\subset\mathbb{R}^3$内找到粒子的概率是
$$P(\text{粒子在Q内})=\int_Q|\psi(x,t)|^2d^3x$$
将薛定谔方程取复共轭后得到
\begin{equation}
    -i\hbar\frac{\partial}{\partial t}\psi^*=-\frac{\hbar^2}{2m}\nabla^2\psi^*+V(x)\psi^*
\end{equation}
将(2.1)式乘以$\psi^*$再减去(2.2)式乘以$\psi$得到
$$i\hbar\psi^*\partial_t\psi+i\hbar\psi\partial_t\psi^*=\frac{\hbar^2}{2m}(-\psi^*\nabla^2\psi+\psi\nabla^2\psi^*)$$
$$\frac{\partial}{\partial t}(\psi^*\psi)+\vec{\nabla}\cdot\left[\frac{i\hbar}{2m}(\psi\vec{\nabla}\psi^*-\psi^*\vec{\nabla}\psi)\right]=0$$
这个式子具有$\partial_t\rho+\vec{\nabla}\cdot\vec{J}=0$的守恒流形式，其中$\rho=\psi^*\psi$为概率密度，$\vec{J}=\frac{i\hbar}{2m}(\psi\vec{\nabla}\psi^*-\psi^*\vec{\nabla}\psi)$为概率流密度。这个式子还告诉我们，全空间中的概率应该随时间不变，也就是说
$$\frac{d}{dt}\int_{\mathbb{R}^3}|\psi(x,t)|^2d^3x=0$$
只要在某个时刻确定归一化条件$\int|\psi|^2d^3x=1$，这个条件随系统演化会自动保持。
\subsection{Hilbert空间}
\subsubsection{基本定义}
Hilbert空间$\mathcal{H}$是一个元素的集合，其上定义有三种运算：数乘、加法和内积。关于数乘和加法，$\mathcal{H}$构成复数域$\mathbb{C}$上的线性空间。关于内积，满足以下性质
\\ (a)内积$(\cdot,\cdot)$是$\mathcal{H}\times\mathcal{H}\to \mathbb{C}$的一个映射，对于$\psi,\phi\in\mathcal{H}$，$(\psi,\phi)\in\mathbb{C}$
\\ (b)反共轭性 $(\psi,\phi)=\overline{(\phi,\psi)}\quad\forall\phi,\psi\in\mathcal{H}$
\\ (c)线性 $(\psi,\alpha\phi_1+\beta\phi_2)=\alpha(\psi,\phi_1)+\beta(\psi,\phi_2)\quad\forall\alpha,\beta\in\mathbb{C},\psi,\phi_1,\phi_2\in\mathcal{H}$
\\ (d)正定性 $(\psi,\psi)\geq 0\quad\forall\psi\in\mathcal{H}$，且$(\psi,\psi)=0$当且仅当$\psi=0$
\\ 类比于$\mathbb{R}^3$中向量的模长，定义$\mathcal{H}$上的范数为
$$\|\psi\|=\sqrt{(\psi,\psi)}$$

注：若追求数学的严格性，还应该添加完备性条件，$\mathcal{H}$中的任意柯西序列都收敛。

\begin{example}
    二维复线性空间$\mathbb{C}^2$，定义内积
    $$(\psi,\phi)=\overline{\psi}_1\phi_1+\overline{\psi}_2\phi_2\quad \psi=\left[\begin{array}{c}
        \psi_1\\
        \psi_2
    \end{array}\right]\quad \phi=\left[\begin{array}{c}
        \phi_1\\
        \phi_2
    \end{array}\right]$$
    则其构成一个Hilbert空间。
\end{example}
\begin{example}
    $\mathbb{R}^n$上的所有平方可积复值函数构成无限维复线性空间$L^2(\mathbb{R}^n)$
    $$L^2(\mathbb{R}^n)=\{f:\left|\int_{\mathbb{R}^n}|f(x)|^2d^nx\right|<\infty\}$$
    在其上定义内积
    $$(f,g)=\int_{\mathbb{R}^n}\overline{f(x)}g(x)d^nx$$
    构成Hilbert空间。

    注：这里的积分应该理解为勒贝格积分而非黎曼积分，勒贝格积分有更好的性质使得可积函数乘以可积函数依旧是可积函数。
\end{example}
\begin{theorem}[Cauchy-Schwarz不等式]
    $\mathcal{H}$是Hilbert空间，则其上的内积满足
    $$|(\phi,\psi)|^2\leq(\phi,\phi)(\psi,\psi)\quad\forall\psi,\phi\in H$$
\end{theorem}
\begin{proof}
    若$\psi=0$，则上式两侧均为0，不等式成立，若$\psi\neq 0$，
    对于$\lambda\in\mathbb{C}$
    $$(\phi+\lambda\psi,\phi+\lambda\psi)=(\phi,\phi)+\lambda(\phi,\psi)+\overline{\lambda}(\psi,\phi)+|\lambda|^2(\psi,\psi)\geq 0$$
    取$$\lambda=i\frac{(\psi,\phi)}{(\psi,\psi)}$$
    则有
    $$(\phi,\phi)-\frac{|(\phi,\psi)|^2}{(\psi,\psi)}\geq 0$$
    证毕
\end{proof}
\subsubsection{泛函与算子}
$\mathcal{H}$是Hilbert空间，$f:\mathcal{H}\to \mathbb{C}$是从$\mathcal{H}$到复数域的线性映射，即满足
$$f(\alpha\psi+\beta\phi)=\alpha f(\psi)+\beta f(\phi)\quad\forall\alpha,\beta\in\mathbb{C},\psi,\phi\in\mathcal{H}$$
则$f$称作$\mathcal{H}$上的线性泛函。若还存在一个实数$A\geq 0$，使得
$$|f(\phi)|\leq A\|\phi\|\quad\forall\phi\in\mathcal{H}$$
则$f$称作有界的。

\begin{theorem}
    $f$是Hilbert空间$\mathcal{H}$上的有界线性泛函，则存在唯一的$\psi_f\in\mathcal{H}$，使得
    $$f(\phi)=(\psi_f,\phi)\quad\forall \phi\in\mathcal{H}$$
\end{theorem}
从这个定理可以看出，$\mathcal{H}$中的向量，和$\mathcal{H}$上的有界线性泛函是一一对应的。换言之，若知道了某个向量$\psi$与$\mathcal{H}$中所有向量的内积，便可以把这个向量唯一确定下来。

$\mathcal{H}$到自身的线性映射称为其上的算子，即满足
$$A(\alpha\psi+\beta\phi)=\alpha A\psi+\beta A\phi\quad \forall\psi,\phi\in \mathcal{H}$$
对于某个算子$A$，定义其共轭算子$A^\dagger$为
$$(A^\dagger\psi,\phi)=(\psi,A\phi)\quad \forall\psi,\phi\in\mathcal{H}$$
这个式子为什么可以确定$A^\dagger$呢？我们看到，把$A^\dagger$作用到任意向量$\psi$上，得到一个新向量$A^\dagger\psi$。为了确定$A^\dagger\psi$，我们把它与$\mathcal{H}$中每个向量$\phi$做内积，内积的值由等式右侧确定。因此根据定理2.2，$A^\dagger\phi$就唯一确定下来了，因此$A^\dagger$也唯一确定了。

若某个算子$H$满足$H^\dagger=H$，我们就称其为自伴算子（厄米算子）。

注：数学上严格地说，这个方式只能定义有界算子的伴算子。对于无界算子，其伴算符还有定义域的问题，自伴和厄米也有所不同。
\subsubsection{算子的特征方程}
$\mathcal{H}$中的一组向量$\{e_\nu\}$，若$\mathcal{H}$中任意向量$\phi$都可写成如下形式
$$\phi=\sum_\nu \phi_\nu e_\nu\quad \phi_\nu\in\mathbb{C}$$
则称$\{e_\nu\}$在$\mathcal{H}$是完备的。换言之，$\mathcal{H}$中任意向量可按$\{e_\nu\}$展开。（若$\{e_\nu\}$还是线性无关的，则展开系数$\phi_\nu$唯一确定）

若基$\{e_\nu\}$还满足
$$(e_\nu,e_\mu)=\delta_{\nu\mu}=\left\{\begin{array}{ll}
    0&,\nu\neq\mu\\
    1&,\nu=\mu
\end{array}\right.$$
则$\{e_\nu\}$称作正交归一基，此时向量$\phi$的展开系数为
$$\phi=\sum_\nu\phi_\nu e_\nu\quad \phi_\nu=(e_\nu, f)$$

注：这里的$\sum_\nu$既可以理解为对可数多个基矢的求和，也可以理解为积分
$$\phi=\int \phi(\nu)e_\nu d\nu$$

对于算子$A$，方程
$$A\psi=\lambda\psi\quad \psi\neq 0,\psi\in\mathcal{H},\lambda\in\mathbb{C}$$
称作$A$的特征方程，$\lambda$称作特征值，$\psi$称作特征向量。
\begin{theorem}
    自伴算子的特征值都是实数。
\end{theorem}
\begin{proof}
    假定$A$自伴，$\lambda,\psi$是其特征值与特征向量。则有
    $$(\psi,A\psi)=\lambda(\psi,\psi)$$
    另一方面，也有
    $$(\psi,A\psi)=\overline{(A\psi,\psi)}=\overline{(\psi,A^*\psi)}=\overline{(\psi,A\psi)}=\overline{\lambda}(\psi,\psi)$$
    两式联立有$\lambda=\overline{\lambda}$，证毕。
\end{proof}
\begin{theorem}
    自伴算子不同特征值对应的特征向量相互正交
\end{theorem}
\begin{proof}
    假定$A$自伴，$\lambda_1,\psi_1$和$\lambda_2,\psi_2$为两组特征值和特征向量，$\lambda_1\neq\lambda_2$，则有
    $$(\psi_1,A\psi_2)=\lambda_2(\psi_1,\psi_2)$$
    另一方面
    $$(\psi_1,A\psi_2)=(A^*\psi_1,\psi_2)=(A\psi_1,\psi_2)=(\lambda_1\psi_1,\psi_2)=\lambda_1(\psi_1,\psi_2)$$
    联立有
    $$(\lambda_1-\lambda_2)(\psi_1,\psi_2)=0$$
    因为$\lambda_1\neq\lambda_2$，所以$(\psi_1,\psi_2)=0$，两者相互正交，证毕。
\end{proof}
\begin{theorem}
    自伴算子$A$的所有特征向量组成$\mathcal{H}$的正交归一基。
\end{theorem}
注：证明只需证特征向量组的完备性。对于有限维Hilbert空间，证明可在《线性代数》中找到（任意厄米矩阵可被对角化）。对于无限维空间，数学上严格的证明极其复杂。事实上，许多自伴算符都不具备常规意义下的特征向量，比如$L^2(\mathbb{R}^3)$中的位置算符$(\hat{X}f)(x)=xf(x)$，其特征值为$x_0$的特征向量为$g_{x_0}(x)=\delta(x-x_0)$，这甚至不是一个常规意义上的“函数”。
\begin{theorem}
    $A,B$是自伴算子，则$[A,B]=AB-BA=0$的充要条件是，存在一组正交归一基$\{e_\nu\}$，使得每个$e_\nu$是$A$和$B$的共同特征向量
    $$Ae_\nu=a_{\nu}e_\nu\quad Be_\nu=b_\nu e_\nu$$
\end{theorem}
注：对于有限维Hilbert空间，证明可在《线性代数》中找到。对于无限维空间，数学上严格的证明极其复杂。
\subsubsection{Dirac符号}
对于$\psi\in\mathcal{H}$，引入左右矢记号
$$\bra{\psi}\quad \ket{\psi}$$
把内积关系记为
$$(\phi,\psi)=\braket{\phi|\psi}$$
定义算符$A$作用于右矢上
$$A\ket{\psi}=\ket{A\psi}$$
我们还希望把$A$作用于左矢，并且$\braket{\psi|A|\phi}$在$A$向左或向右作用时相同
$$(\bra{\psi}A)\ket{\psi}=\bra{\psi}(A\ket{\phi})=(\psi,A\phi)$$
因此可定义
$$\bra{\psi}A=\bra{A^\dagger\psi}$$
向量$\psi$按照正交归一基$\{e_\nu\}$的展开现在变成
$$\ket{\psi}=\sum_\nu\ket{e_\nu}\braket{e_\nu|\psi}$$
$$\bra{\psi}=\sum_\nu\braket{\psi|e_\nu}\bra{e_\nu}$$
因此可记
$$\sum_\nu\ket{e_\nu}\bra{e_\nu}=1$$
或
$$\int \ket{e_\nu}d\nu\bra{e_\nu}=1$$

\subsubsection{表象与表象变换}

一组正交归一基$\{\ket{\nu}\}$的好处是可以帮助简化计算。
$$\ket{\psi}=\sum_\alpha\psi_\alpha\ket{\alpha}$$
$$\ket{\phi}=\sum_\alpha\phi_\alpha\ket{\alpha}$$
向量间加法数乘运算变为
$$\ket{\lambda\psi+\phi}=\sum_\alpha(\lambda\psi_\alpha+\phi_\alpha)\ket{\alpha}$$
$$(\lambda\psi+\phi)_\alpha=\lambda\psi_\alpha+\phi_\alpha$$
内积运算变为
$$\braket{\psi|\phi}=\sum_{\alpha\beta}\psi^*_\alpha\phi_\beta\braket{\alpha|\beta}=\sum_{\alpha\beta}\psi^*_\alpha\phi_\beta\delta_{\alpha\beta}=\sum_\alpha\psi^*_\alpha\phi_\alpha$$
算符运算变为
$$A\ket{\psi}=\sum_\alpha \psi_\alpha A\ket{\alpha}=\sum_{\alpha\beta}\braket{\beta|A|\alpha}\psi_\alpha\ket{\beta}$$
$$(A\psi)_\alpha=\sum_\beta A_{\alpha\beta}\psi_\beta\quad A_{\alpha\beta}=\braket{\alpha|A|\beta}$$
$$A=\sum_{\alpha\beta}\ket{\alpha}\braket{\alpha|A|\beta}\bra{\beta}=\sum_{\alpha\beta}\ket{\alpha}A_{\alpha\beta}\bra{\beta}$$
如此便把陌生的Hilbert空间向量，转化成了我们熟悉的矩阵运算。右矢$\ket{}$由列向量$\psi_\alpha$描述，矩阵有方阵$A_{\alpha\beta}$描述。这便称作$\{\ket{\alpha}\}$表象。

假设有另外一套正交归一基$\{\ket{a}\}$，在这个表象下，$\psi$和$A$表示为
$$\psi_a=\braket{a|\psi}=\sum_\alpha\braket{a|\alpha}\braket{\alpha|\psi}=\sum_\alpha U_{a\alpha}\psi_\alpha$$
$$A_{ab}=\braket{a|A|b}=\sum_{\alpha\beta}\braket{a|\alpha}\braket{\alpha|A|\beta}\braket{\beta|b}=\sum_{\alpha\beta}U_{a\alpha}U_{b\beta}^*A_{\alpha\beta}$$
$U_{a\alpha}=\braket{a|\alpha}$是两个表象间的变换矩阵，其满足
$$\sum_{\alpha}U_{a\alpha}U^\dagger_{\alpha b}=\sum_{\alpha}U_{a\alpha}U^*_{b\alpha}=\sum_{\alpha}\braket{a|\alpha}\braket{\alpha|b}=\braket{a|b}=\delta_{ab}$$
因此$U_{a\alpha}$是幺正矩阵。
\begin{example}
    在$L^2(\mathbb{R}^3)$中，有坐标算符$\hat{X}$与动量算符$\hat{P}$，分别有一套本征矢$\ket{x},\ket{p}$（取$\hbar=1$）
    $$\ket{x_0}(x)=\delta(x-x_0)\quad \ket{p}(x)=\frac{1}{(2\pi)^{3/2}}\e^{\i \vec{p}\cdot\vec{x}}$$
    $$\braket{x_1|x_2}=\delta(x_1-x_2)\quad \braket{p_1|p_2}=\delta(p_1-p_2)$$
    $$\braket{x|\psi}=\psi(x)\quad \braket{p|\psi}=\frac{1}{(2\pi)^{3/2}}\int \e^{\i \vec{p}\cdot\vec{x}}\psi(x)d^3x$$
    $$\braket{x_2|\hat{X}|x_1}=x_1\braket{x_2|x_1}=x_1\delta(x_1-x_2)$$
    $$\braket{x_2|\hat{P}|x_1}=\int\braket{x_2|p_2}\braket{p_2|\hat{P}|p_1}\braket{p_1|x_1}d^3p_1d^3p_2=\frac{1}{(2\pi)^{3}}\int \e^{\i p_1x_1-\i p_2x_2}p_1\delta(p_1-p_2)d^3p_1d^3p_2$$
    $$=\frac{1}{(2\pi)^3}\int p\e^{\i p(x_1-x_2)}d^3p=-\i\nabla\delta(x_1-x_2)$$
    注：这些计算在数学上是不严格的，$\ket{x}$和$\ket{p}$都不属于$L^2(\mathbb{R}^3)$，两个向量内积得出$\delta$函数也不合理。数学上严格的做法是以广义本征矢来理解这些向量。

\end{example}

\subsection{量子力学基本假设}

\begin{postulate}[量子态假设]
    物理系统由一个Hilbert空间$\mathcal{H}$表示。每一个时刻，系统的状态由$\mathcal{H}$中的归一化向量$\psi$描述。相差一个相因子$\e^{\i\delta}$的向量描述同一个系统状态。也就是说，系统的状态与$\mathcal{H}$中的射线一一对应。
\end{postulate}
\begin{postulate}[算符假设]
    系统中每一个可观测力学量$Q$，由$\mathcal{H}$中的自伴算子$\hat{Q}$描述。力学量之间的运算关系由算符间的运算实现。
\end{postulate}
\begin{example}
    一个在3维空间中运动的粒子，描述这个系统的Hilbert空间是$L^2(\mathbb{R}^3)$。粒子的位置由位置算符$\hat{X}_i$描述，动量由动量算符$\hat{P}_i$描述
    $$(\hat{X}_if)(x)=x_if(x)\quad (\hat{P}_if)(x)=-\i\hbar\dfrac{\partial f}{\partial x_i}(x)$$
    角动量与位置、动量的关系为
    $$\vec{L}=\vec{x}\times \vec{p}$$
    这个关系由算符的运算实现
    $$\hat{L}_i=\sum_{jk}\epsilon_{ijk}\hat{X}_j\hat{P}_k$$
    动能与动量的关系为
    $$T=\frac{p^2}{2m}\quad \hat{T}=\sum_i\frac{\hat{P}_i^2}{2m}\quad (\hat{T}f)(x)=-\frac{\hbar^2}{2m}\nabla^2 f(x)$$
\end{example}
\begin{postulate}[测量假设]
    对某个可观测力学量$Q$进行测量，测量所得结果是$\hat{Q}$的某个特征值$q$。测量过程会导致系统状态$\psi$变为$q$对应的特征向量$\psi_q$，得到$q$值的概率是
    $$P\left[Q=q|\psi\right]=|(\psi,\psi_q)|^2$$
\end{postulate}
\begin{theorem}
    在状态$\psi$下测量可观测力学量$Q$，得到结果的期望值为
    $$\mathrm{E}_\psi[Q]=\braket{\psi|\hat{Q}|\psi}$$
\end{theorem}
\begin{proof}
    $$\mathrm{E}[Q]=\sum_qP[Q=q]q=\sum_q \braket{\psi|\psi_q}q\braket{\psi_q|\psi}$$
    $$=\sum_q\braket{\psi|A|\psi_q}\braket{\psi_q|\psi}=\braket{\psi|A|\psi}$$
    最后一步利用了自伴算符特征向量的完备性条件
\end{proof}
注：这里对测量结果的期望，应该理解为制备多个系统$\psi$，对每个系统测量得到一个结果，再取平均。而非对某个系统进行连续的多次测量。
\begin{theorem}[不确定关系]
    两个可观测力学量$A,B$，在状态$\psi$下测量的不确定度定义为
    $$\Delta_\psi A=\sqrt{\mathrm{E}_\psi\left[(A-\mathrm{E}_\psi[A])^2\right]}$$
    有
    $$\Delta_\psi A\Delta_\psi B\geq\frac{1}{2}\left|\braket{\psi|[\hat{A},\hat{B}]|\psi}\right|$$
\end{theorem}
\begin{proof}
    首先证明，$\braket{\psi|[\hat{A},\hat{B}]|\psi}$是纯虚数。
    $$\overline{(\psi,[\hat{A},\hat{B}]\psi)}=([\hat{A},\hat{B}]\psi,\psi)=(\psi,[\hat{A},\hat{B}]^\dagger\psi)=-(\psi,[\hat{A},\hat{B}]\psi)$$
    对于任意$\xi\in\mathbb{R}$
    $$I(\xi)=\|(\xi\hat{A}+i\hat{B})\psi\|^2=(\xi\hat{A}\psi+i\hat{B}\psi,\xi\hat{A}\psi+i\hat{B}\psi)\geq 0$$
    换一种写法
    $$I(\xi)=\xi^2(\hat{A}\psi,\hat{A}\psi)+(\hat{B}\psi,\hat{B}\psi)+i\xi(\hat{A}\psi,\hat{B}\psi)-i\xi(\hat{B}\psi,\hat{A}\psi)$$
    $$=\xi^2\braket{\psi|\hat{A}^2|\psi}+\braket{\psi|\hat{B}^2|\psi}+\xi\braket{\psi|i(\hat{A}\hat{B}-\hat{B}\hat{A})|\psi}\geq 0$$
    因此这个关于$\xi$的二次函数的判别式$\leq 0$
    $$\braket{\psi|i[\hat{A},\hat{B}]|\psi}^2\leq 4\braket{\psi|\hat{A}^2|\psi}\braket{\psi|\hat{B}^2|\psi}$$
    再将$\hat{A}$替换为$\hat{A}-\braket{\psi|\hat{A}|\psi}$，将$\hat{B}$替换为$\hat{B}-\braket{\psi|\hat{B}|\psi}$便证毕。
\end{proof}

在前面的部分，我们看到，相互对易的自伴算符有共同的本征态，并且这些本征态构成$\mathcal{H}$的正交归一基。这有什么用呢？注意到，对于某个可观测力学量$F_0$，其每个本征值对应的本征向量空间可能是简并的，即一个本征值对应多个线性无关本征向量。为了在简并子空间中再挑选一组向量作为正交归一基，我们可以再找一个与$F_0$对易的力学量$F_1$，找$(F_0,F_1)$的共同本征态，看是否能消除简并，如果还不可以，就再找$F_3,...,F_n,...$，一直到消除全部简并。
\begin{postulate}[时间演化假设]
    量子态$\psi$随时间的演化满足薛定谔方程
    $$\i\hbar\frac{\partial\psi}{\partial t}=\hat{H}\psi$$
    $\hat{H}$是系统哈密顿量对应的哈密顿算符。
\end{postulate}
\begin{postulate}[正则量子化方案]
    系统的坐标和动量间满足正则对易关系$$[\hat{X}_i,\hat{P}_j]=\i\hbar\delta_{ij}$$
\end{postulate}

这里我们说到到“量子化”，那什么是量子化？量子化就是构造一个Hilbert空间，指定可观测力学量到算符的对应关系，使上述量子力学的假设可以在这个Hilbert空间上自洽地成立。比如说，物理系统是一个粒子在空间中运动，我们找到的Hilbert空间就是$L^2(\mathbb{R}^3)$，再指定$\hat{X}=x\quad\hat{P}=-\i\hbar\nabla$，这样正则对易关系成立，完成了量子化。
\subsection{例：代数法解谐振子}

\subsection{绘景}
我们到现在一直认为，基本力学量$X,P$是$\mathcal{H}$中的算符，不随时间变化，而系统状态$\psi$随时间的演化
$$\i\hbar\frac{\partial\psi}{\partial t}=\hat{H}\psi=H(\hat{X},\hat{P},t)\psi(t)$$
我们按照下式定义一个算符$U(t,t_0)$，称为时间演化算符
$$U(t_0,t_0)=I\quad\i\hbar\frac{\partial}{\partial t}U(t,t_0)=H(\hat{X},\hat{P},t)U(t,t_0)$$
\begin{theorem}
    $U$满足以下性质：\\
    (a) $U(t,t_0)^\dagger U(t,t_0)=I$\\
    (b) $U(t_2,t_1)U(t_1,t_0)=U(t_2,t_0)$\\
    (c) 对于满足薛定谔方程的$\psi(t)$，$\psi(t)=U(t,t_0)\psi(t_0)$
\end{theorem}
\begin{proof}
    先证(a)，将$U$满足的方程取共轭后有
    $$\i\hbar\frac{\partial}{\partial t}U(t,t_0)^\dagger=-U(t,t_0)^\dagger\hat{H}(t)$$
    $$\i\hbar\frac{\partial}{\partial t}\left[U(t,t_0)^\dagger U(t,t_0)\right]=U(t,t_0)^\dagger \hat{H}(t)U(t,t_0)-U(t,t_0)^\dagger \hat{H}(t)U(t,t_0)=0$$
    $$U(t_0,t_0)^\dagger U(t_0,t_0)=I$$
    因此$U(t,t_0)^\dagger U(t,t_0)=I$

    再证(b)，记等号左侧为$L_{t_1,t_0}(t_2)$，等号右侧为$R_{t_1,t_0}(t_2)$，满足
    $$\i\hbar\frac{\partial}{\partial t_2}L_{t_1,t_0}(t_2)=\hat{H}(t_2)L_{t_1,t_0}(t_2)$$
    $$\i\hbar\frac{\partial}{\partial t_2}R_{t_1,t_0}(t_2)=\hat{H}(t_2)R_{t_1,t_0}(t_2)$$
    $$ L_{t_1,t_0}(t_1)=R_{t_1,t_0}(t_1)=U(t_1,t_0)$$
    $L$和$R$满足同一方程与初值条件，因此$L=R$。(c)显然易见。
\end{proof}

我们称这样的对系统的描述，态矢量含时算符不含时为薛定谔绘景，相应的态矢量和算符记为$\psi^S(t),A^S$。我们再定义一套新的表示法。先取定某个时刻$t_0$，用$t_0$时的态矢量代表整个演化过程。
$$\psi^H=U(t,t_0)^\dagger \psi^S(t)=\psi^S(t_0)$$
让算符含时演化
$$A^H(t)=U(t,t_0)^\dagger A^S(X^S,P^S,t)U(t,t_0)$$

此时有
$$\braket{\psi^H|A^H(t)|\phi}=\braket{U(t,t_0)^\dagger\psi^S(t)|U(t,t_0)^\dagger A^S(t)U(t,t_0)U(t,t_0)^\dagger|\phi^S(t)}=\braket{\psi^S(t)|A^S(t)|\phi^S(t)}$$
$$\i\hbar\frac{\d}{\d t}A^H(t)=-U^\dagger H^S A^SU+\i\hbar U^\dagger \frac{\partial}{\partial t}A^S(X^S,P^S,t)U+U^\dagger A^S H^SU$$
$$=-U^\dagger H^SUU^\dagger A^SU+U^\dagger A^SUU^\dagger H^SU+U^\dagger\frac{\partial}{\partial t}A^S U$$
$$=[A^H,H^H]+\i\hbar\left(\frac{\partial A}{\partial t}\right)^H$$
这样的对系统的描绘称作海森堡绘景。

物理量$A$的期望值不随时间变化的充要条件是
$$\frac{\d}{\d t}A^H(t)=0$$
若$A$仅是$X,P$的函数而不显含时间，则$A$是守恒量的充要条件是
$$[A^S,H^S]=0$$

\subsection{量子力学的路径积分形式}
考虑一维系统，采用薛定谔表象，我们试图计算
$$\braket{r'|U(t',t_0)|r_0}\quad t_0<t'$$
将$[t_0,t']$均匀划做$N$份，$t_0,t_1,...,t_N=t'$，将$U(t',t_0)$拆为$N$个$U(t+\Delta t,t)$的乘积
$$U(t',t_0)=U(t',t_{N-1})U(t_{N-1},t_{N-2}),...,U(t_1,t_0)$$
再插入多个完备性关系$\int\ket{r}\bra{r}dr=1$
$$\braket{r'|U(t',t_0)|r_0}=\int\braket{r'|U(t',t_{N-1})|r(t_{N-1})}\braket{r(t_N-1)|U(t_{N-1},t_{N-2})}......$$
$$\braket{r(t_1)|U(t_1,t_0)|r_0}dr(t_{N-1})...dr(t_1)$$
取$\Delta t\to 0$的极限，可以有
$$U(t+\Delta,t)\approx 1-\i\hbar H(t)\Delta t$$
$$\braket{r(t_{i+1})|U(t_{i+1},t_i)|r(t_i)}=\braket{r(t_{i+1})|1-\frac{\i}{\hbar} H(P,X,t)\Delta t|r(t_i)}$$
假设$H$中所有的$P$都写在所有$X$左侧，再插入一个动量本征态完备条件$\int\ket{p(t_{i+1})}\bra{p(t_{i+1})}dp(t_{i+1})=1$，有
$$=\int\braket{r(t_{i+1})|p(t_{i+1})}\braket{p(t_{i+1})|1-\frac{\i}{\hbar} H(P,X,t)\Delta t|r(t_i)}dp(t_{i+1})$$
$$\approx \int\frac{1}{2\pi\hbar}\exp\left\{-\frac{\i}{\hbar} p(t_{i+1})\left[r(t_{i+1})-r(t_i)\right]-\frac{\i}{\hbar}H\left[p(t_{i+1}),r(t_i),t_i\right]\right\}dp(t_{i+1})$$
因此有
$$\braket{r'|U(t',t_0)|r_0}=\lim_{N\to\infty}\int_{r(t_0)=r_0;r(t_N)=r'}\exp\left\{-\frac{\i}{\hbar}\sum_{i=0}^{N-1}\left\{p(t_{i+1})\left[r(t_{i+1})-r(t_i)\right]+H\left[p(t_{i+1}),r(t_i),t_i\right]\Delta t\right\}\right\}$$
$$dr(t_1)dr(t_2)...dr(t_{N-1})\frac{dp(t_1)}{2\pi\hbar}\frac{dp(t_2)}{2\pi\hbar}...\frac{dp(t_{N})}{2\pi\hbar}$$
若$H$是动量的二次型，$H=\frac{p^2}{2m}+V(r,t)$，则可以完成对$p$的积分，利用关系
$$\int dx \exp\left[-\i bx-\i ax^2\right]=\sqrt{\frac{\pi}{\i a}}\exp\left[\i\frac{b^2}{4a}\right]$$
$$\int \exp\left\{-\frac{\i}{\hbar}\left\{p(t_{i+1})\left[r(t_{i+1})-r(t_i)\right]+H\left[p(t_{i+1}),r(t_i),t_i\right]\Delta t\right\}\right\}\frac{dp(t_{i+1})}{2\pi\hbar}$$
$$=\sqrt{\frac{m}{2\pi \i\hbar\Delta t}}\exp\left\{\frac{\i}{\hbar}\left[\frac{1}{2}m\left(\frac{r(t_{i+1})-r(t_i)}{\Delta t}\right)^2-V(r(t_i),t_i)\right]\Delta t\right\}$$
$$\braket{r'|U(t',t_0)|r_0}=\lim_{N\to\infty}\left(\frac{m}{2\pi \i\hbar\Delta t}\right)^{\frac{N}{2}}$$
$$\int \exp\left\{\frac{\i}{\hbar}\sum_{i=0}^{N-1}\left[\frac{1}{2}m\left(\frac{r(t_{i+1})-r(t_i)}{\Delta t}\right)^2-V(r(t_i),t_i)\right]\Delta t\right\}\prod_{i=1}^{N-1}dr(t_i)$$
可记作
$$\braket{r'|U(t',t_0)|r_0}\propto \int_{r(t_0)=r_0;r(t')=r'}[\mathcal{D}r]\exp\left[\frac{\i}{\hbar} S[r(t)]\right]$$

\section{对称性与守恒量}
\subsection{一般性论述}
\subsubsection{保持跃迁概率的变换}
对称性在物理学中有极其重要的地位。若系统在某个变换下性质保持不变，我们称系统具有对于这个变换的对称性。所以我们先来研究变换。

长久以来，对于变换有两种理解方式。一种是主动的观点，认为变换是系统自己在变动；另一种是被动的观点，认为系统自己没有发生变化，变动的是描述系统的方式。以平移变换为例，主动变换的观点是整个系统（粒子、场）向前平移了一段距离$a$，被动变换的观点是系统仍在原地不变，只是观测系统的方式（实验室、仪器、科研人员）整体向后平移了$a$。两种观点在数学上的描述是等同的，我们采用主动的观点。

回到量子力学，在这里系统的状态由希尔伯特空间中的单位向量$\psi$描述，物理量由算符描述。考虑某个变换$\alpha$，按照主动变换的观点，它应该把系统状态从$\psi$变为$\psi_\alpha$，但不改变描述系统的方式，即不改变物理量的算符。

变换有了，对称性又是什么呢？应该是系统的“性质”在变换下不变。所以我们来考察量子力学中系统的哪些性质应当不变。

最简单最一般的性质，应该是系统状态之间的跃迁概率$P_{\psi\to\phi}|(\psi,\phi)|^2$，它在变换$\alpha$下不变的数学描述是
$$|(\psi_\alpha,\phi_\alpha)|^2=|(\psi,\phi)|^2\quad\forall\psi,\phi$$
Wigner定理告诉我们，保持跃迁概率不变的$\alpha$，可以用幺正算子或反幺正算子表示。
\begin{theorem}[Wigner]
    $\mathcal{H}$是一个Hilbert空间，$\mathcal{P}$是其上所有射线组成的集合，$\alpha:\psi\to\psi_\alpha$是$\mathcal{P}$到自身的一一对应的映射（$\psi$和$\psi_\alpha$是用来代表它们所在射线的向量）。若$\alpha$满足
    $$|(\psi_\alpha,\phi_\alpha)|^2=|(\psi,\phi)|^2\quad \forall\psi,\phi$$
    则$\mathcal{H}$上存在一个幺正算子（或反幺正算子）$U_\alpha$，使得
    $$\psi_\alpha=U_\alpha\psi$$
    （上式的确切意义是，$U_\alpha\psi$在$\psi_\alpha$代表的射线上）
\end{theorem}
\subsubsection{变换群}
对称性更美好的地方，在于对称性变换可以构成群。定义两个变换$\alpha,\beta$的乘积$(\alpha\beta)$为先执行$\beta$，再执行$\alpha$
$$\psi_{(\alpha\beta)}=\left(\psi_{\beta}\right)_{\alpha}$$
若$\alpha,\beta$都保持跃迁概率不变，则$(\alpha\beta)$也保持跃迁概率不变，由Wigner定理，它们都可以用（反）幺正算符刻画
$$\psi_{(\alpha\beta)}=U_{(\alpha\beta)}\psi\quad \psi_\alpha=U_\alpha\psi\quad \psi_\beta=U_\beta\psi$$
再代入乘积$(\alpha\beta)$的定义有
$$U_{(\alpha\beta)}\psi=\psi_{(\alpha\beta)}=(\psi_\beta)_\alpha=U_\alpha\psi_\beta=U_\alpha U_\beta\psi$$
从上式貌似可以直接推出
$$U_{(\alpha\beta)}=U_\alpha U_\beta$$
但实际要困难许多，因为这些“相等”都是射线的相等，可以随时随地差一个相因子
$$U_{(\alpha\beta)}\psi=e^{i\delta(\alpha,\beta,\psi)}U_\alpha U_\beta\psi$$
幸运的是，对于量子力学中绝大多数情况，都可以通过重定义$U$，消去这些相因子，使得
$$U_{(\alpha\beta)}=U_\alpha U_\beta\quad\forall\alpha,\beta\in g$$
这样，我们得到了变换群$g$在Hilbert空间上的幺正表示$U$。（由于反幺正算子的乘积是幺正算子，因此反幺正算子不能构成群）

可以更进一步，得到量子力学版本的Noether定理
\begin{theorem}\label{Noether}
    设$U(\lambda)$是Hilbert空间$\mathcal{H}$上的幺正算子群，$\lambda\in\mathbb{R}$是群参数，且群乘法为
    $$U(\lambda)U(\kappa)=U(\lambda+\kappa)\quad \lambda,\kappa\in\mathbb{R}$$
    则$\mathcal{H}$上存在唯一的自伴算子$A$，使得
    $$A=i\left.\frac{\d}{\d\lambda}U(\lambda)\right|_{\lambda=0}\quad U(\lambda)=\exp(-i\lambda A)\quad\forall\lambda\in\mathbb{R}$$
    $A$称为生成元。
\end{theorem}
\begin{proof}
    这个定理有一个不严格的证明。首先因为$U(0)U(0)=U(0)$，所以
    $$U(0)=I$$
    令
    $$A=i\left.\frac{\d}{\d\lambda}U(\lambda)\right|_{\lambda=0}\quad V(\lambda)=\exp(-i\lambda A)$$
    对$\lambda$求导有
    $$\frac{\d}{\d\lambda}V(\lambda)=-iA\exp(-i\lambda A)=-iAV(\lambda)$$
    另一方面
    $$\frac{\d}{\d\lambda}U(\lambda)=\lim_{\delta\to 0}\frac{U(\lambda+\delta)-U(\lambda)}{\delta}=\lim_{\delta\to 0}\frac{U(\delta)-U(0)}{\delta}U(\lambda)=-iAU(\lambda)$$
    $U(\lambda)$与$V(\lambda)$满足相同的边界条件和微分方程，因此$U\equiv V$。
\end{proof}
只是这个定理还不能直接地看出其反映对称性与守恒量间的关系，我们再接着分析。假设现在有一个对称变换群$g$，其有一个单参数幺正算子表示
$$U(\lambda)=e^{-i\lambda A}$$
再假设系统演化方程——薛定谔方程在这个变换群下不变，换言之，$\psi(t)$和$U(\lambda)\psi(t)$均满足薛定谔方程
$$i\frac{\partial}{\partial t}\psi(t)=H(t)\psi(t)\quad i\frac{\partial}{\partial t}U(\lambda)\psi(t)=H(t)U(\lambda)\psi(t)$$
很清楚地看到，上式对于任意$\psi(t)$均成立的条件是
$$H(t)U(\lambda)=U(\lambda)H(t)\quad [U(\lambda),H(t)]\quad \forall t,\lambda$$
又因为$A=i\frac{\d}{\d\lambda}U(\lambda)$，有
$$[A,H(t)]=i\frac{\d}{\d \lambda}[U(\lambda),H(t)]=0$$
对于哈密顿量不含时的情况，$A$便是一个守恒量。也就是说，一个保持系统演化方程不变的对称性，对应一个守恒量。

再来看看物理量在对称性变换$\alpha$下如何变动。考虑物理量$F$，对称性变换$\alpha$和其对应的幺正算子$U_{\alpha}$。则$F$在变换后的期望值是
$$(\psi_\alpha,F\psi_\alpha)=(U_\alpha\psi,FU_\alpha\psi)=\braket{\psi|U^{-1}_\alpha FU_\alpha|\psi}$$
若物理量$F$有关于这个变换的对称性，在变换下不变，它的期望值应当不变，换言之
$$U^{-1}_\alpha FU_\alpha=F\quad [U_\alpha,F]=0$$
若$U_\alpha$是按单参数群的方式呈现，$U(\alpha)=e^{i\alpha A}$，则上述关系还可以写为
$$[A,F]=0$$
\subsection{空间平移}
若空间平移是某个系统的对称性，可以保持其跃迁概率不变，则在Hilbert空间中存在空间平移变换群的幺正表示$U(\vec{a})$，$\vec{a}\in\mathbb{R}^3$，其有群乘法
$$U(\vec{a})U(\vec{b})=U(\vec{a}+\vec{b})$$
可以把每个平移拆成三个方向平移之和
$$U(\vec{a})=U(a_1\hat{e}_1+a_2\hat{e}_2+a_3\hat{e}_3)=U(a_1\hat{e}_1)U(a_2\hat{e}_2)U(a_3\hat{e}_3)$$
每个单独方向的$U(a_i\hat{e}_i)$都满足定理\ref{Noether}，因此可有一系列自伴算子$P_i$，使得
$$U(a_i\hat{e}_i)=e^{-\i P_ia_i}\quad i=1,2,3$$
因为不同的$U(\vec{a})$之间是对易的，所以其生成元也是对易的
$$[P_i,P_j]=0\quad i,j=1,2,3$$
因此可以把$U(\vec{a})$写为
$$U(\vec{a})=U(a_1\hat{e}_1)U(a_2\hat{e}_2)U(a_3\hat{e}_3)=e^{-\i\vec{a}\cdot\vec{P}}$$

例如，在$\mathrm{L}^2(\mathbb{R}^3)$中，定义平移变换
$$\left[U(a)\psi\right](\vec{x})=\psi(\vec{x}-\vec{a})$$
其图像正是把整个函数$\psi$平移了$\vec{a}$，其生成元为
$$\left[P_i\psi\right](\vec{x})=i\frac{\partial}{\partial a_i}\left[U(a)\psi\right](x)=i\left.\frac{\partial}{\partial a_i}\psi(\vec{x}-\vec{a})\right|_{\vec{a}=0}=-i\frac{\partial}{\partial x_i}\psi(\vec{a})$$
可以看到正是动量算符，事实上，我们可以定义动量算符为平移变换的生成元。
\subsection{时间演化}
这一小节，我们从对称性的角度“推导”薛定谔方程。考虑对系统的一个变换$\alpha(t,t_0)$，其意义是，对于$t_0$时的系统状态，让其自然演化至$t$时刻。假定这个变换保持跃迁概率不变，则其由幺正或反幺正的算子$U(t,t_0)$刻画。

我们再为$U(t,t_0)$做一些合理的假定
$$U(t,t_0)U(t_0,t_1)=U(t,t_1)$$
$$U(t_0,t_0)=I$$
$$U(t,t_0)\psi\ \text{关于}t\text{连续}$$
由上述知$U(t,t_0)$不能是反幺正算子。（$I$不可能连续地变为反幺正算子）定义
$$H(t,t_0)=i\frac{\partial U(t,t_0)}{\partial t} U(t_0,t)=-iU(t,t_0)\frac{\partial U(t_0,t)}{\partial t}$$
$H(t)=H(t,t_0)$是与$t_0$无关的自伴算子，根据下式可以看出
$$H(t,t_0)=i\frac{\partial U(t,t_0)}{\partial t}U(t_0,t)=i\frac{\partial U(t,t_1)U(t_1,t_0)}{\partial t}U(t_0,t_1)U(t_1,t)=i\frac{\partial U(t,t_1)}{\partial t}U(t_1,t)=H(t,t_1)$$
$$H(t)^\dagger=-i U(t_0,t)^\dagger\frac{\partial U(t,t_0)^\dagger}{\partial t}=-iU(t,t_0)\frac{\partial U(t_0,t)}{\partial t}=H(t)$$
并且
$$i\frac{\partial U(t,t_0)}{\partial t}=H(t)U(t,t_0)$$
$H(t)$就是薛定谔方程中的哈密顿量。

再来考虑薛定谔方程在时间平移下的变换，薛定谔方程在时间平移下不变，是指对于任意满足薛定谔方程的$\psi(t)$，平移一段时间$\Delta $后的态$\psi_\Delta(t)=\psi(t-\Delta)$也满足薛定谔方程
$$i\frac{\partial}{\partial t}\psi_\Delta (t)=H(t)\psi_\Delta(t)$$
显然$\psi_\Delta(t)$满足
$$i\frac{\partial}{\partial t}\psi_\Delta(t)=H(t-\Delta)\psi_\Delta(t)$$
两式若相等，只有
$$H(t)=H(t-\Delta)$$
即薛定谔方程在时间平移下不变$\Leftrightarrow$体系哈密顿量不含时$\Leftrightarrow$哈密顿量是守恒量

\section{角动量理论}
\subsection{SO(3)}
三维实正交矩阵群O(3)的定义是
$$\mathrm{O}(3)=\{A\in\mathbb{R}^{3\times 3}:A^TA=I\}$$
由$A^TA=I$，两侧取行列式有
$$(\det A)^2=1\quad \det A=\pm 1$$
由此可知，O(3)群至少有两个分叉，一个分叉为$\det A=1$，另一个为$\det A=-1$，两个分叉间是不连通的。定义
$$\mathrm{SO}(3)=\{A\in\mathrm{O}(3):\det A=1\}$$
\begin{theorem}
    $\mathrm{SO}(3)$中任意群元$A$都可以写成绕某个轴$k$的转动$C_k(\theta)$。即对于任意$a\in\mathbb{R}^3$，$Aa$相当于把$a$绕$k$转动$\theta$角。
\end{theorem}
\begin{proof}
    先证明存在$k\neq 0$，使得$Ak=k$，即证明$\det(A-I)=0$
    $$\det(A-I)=\det(A-I)^T=\det(A^T-I)=\det(A^{-1}-I)=\det A^{-1}\det(I-A)=-\det(A-I)$$
    找到$k$后，再选取两个与$k$垂直的向量，使其构成右手基矢$(i,j,k)$，在该基矢下，$A$应该写为
    $$A=\left[\begin{array}{ccc}
        a_{11}&a_{12}&0\\
        a_{21}&a_{22}&0\\
        0&0&1
    \end{array}\right]$$
    再利用$A^TA=1$的条件，可知$A$唯一可能的形式是
    $$A=\left[\begin{array}{ccc}
        \cos\theta&-\sin\theta&0\\
        \sin\theta&\cos\theta&0\\
        0&0&1
    \end{array}\right]$$
    这正是绕$k$轴旋转$\theta$的矩阵。
\end{proof}
由于SO(3)的任意群元$A$可以写为$A=C_k(\theta)$，保持$k$不变，让$\theta$连续地变化至0，便可把$A$连续的变为$I$，因此我们还证明了
\begin{theorem}
    $\mathrm{SO}(3)$是连通的。
\end{theorem}
接下来，我们看看$C_k(\theta)$在原本的基下的矩阵是什么。经过简单的坐标变换，应当有
$$C_k(\theta)=\left[\begin{array}{ccc}
    i&j&k
\end{array}\right]\left[\begin{array}{ccc}
        \cos\theta&-\sin\theta&0\\
        \sin\theta&\cos\theta&0\\
        0&0&1
    \end{array}\right]\left[\begin{array}{c}
    i^T\\
    j^T\\
    k^T
\end{array}\right]$$
再算一个$C_k(\theta)$对$\theta$的导数
$$\frac{\d}{\d\theta}C_k(\theta)=\left[\begin{array}{ccc}
    i&j&k
\end{array}\right]\left[\begin{array}{ccc}
        0&-1&0\\
        1&0&0\\
        0&0&0
    \end{array}\right]\left[\begin{array}{ccc}
        \cos\theta&-\sin\theta&0\\
        \sin\theta&\cos\theta&0\\
        0&0&1
    \end{array}\right]\left[\begin{array}{c}
    i^T\\
    j^T\\
    k^T
\end{array}\right]$$
$$=\left[\begin{array}{ccc}
    i&j&k
\end{array}\right]\left[\begin{array}{ccc}
        0&-1&0\\
        1&0&0\\
        0&0&0
    \end{array}\right]\left[\begin{array}{c}
    i^T\\
    j^T\\
    k^T
\end{array}\right]C_k(\theta)=(-ij^T+ji^T)C_k(\theta)$$
使用条件$k=i\times j$，有
$$(-ij^T+ji^T)=\left[\begin{array}{ccc}
    0&-i_1j_2+i_2j_1&-i_1j_3+i_3j_1\\
    -i_2j_1+i_1j_2&0&-i_2j_3+i_3j_2\\
    -i_3j_1+i_1j_3&-i_3j_2+i_2j_3&0
\end{array}\right]=\left[\begin{array}{ccc}
    0&-k_3&k_2\\
    k_3&0&-k_1\\
    -k_2&k_1&0
\end{array}\right]$$
$$\frac{\d}{\d\theta}C_k(\theta)=\i(\vec{k}\cdot\vec{J})C_k(\theta)$$
$$J_1=\left[\begin{array}{ccc}
    0&0&0\\
    0&0&\i\\
    0&-\i&0
\end{array}\right]\quad J_2=\left[\begin{array}{ccc}
    0&0&-\i\\
    0&0&0\\
    \i&0&0
\end{array}\right]\quad J_3=\left[\begin{array}{ccc}
    0&\i&0\\
    -\i&0&0\\
    0&0&0
\end{array}\right]$$
\begin{theorem}
    $$C_k(\theta)=\exp\left[\i\theta(\vec{k}\cdot\vec{J})\right]$$
\end{theorem}
\begin{proof}
    $C_k(\theta)$和$A(\theta)=\exp\left[\i\theta(\vec{k}\cdot\vec{J})\right]$都满足同一个微分方程与边界条件
    $$\frac{\d}{\d\theta}A(\theta)=\i(\vec{k}\cdot\vec{J})A(\theta)\quad A(0)=I$$
    证毕
\end{proof}
$\vec{J}$称作SO(3)的生成元，其满足对易关系
$$[J_i,J_j]=\i\epsilon_{ijk}J_k\quad i,j=1,2,3$$
\subsection{SU(2)}
二维复正交幺模矩阵群$\mathrm{SU}(2)$的定义是
$$\mathrm{SU}(2)=\{A\in\mathbb{C}^{2\times 2}:AA^\dagger=I,\det A=1\}$$
\begin{theorem}\label{SU2_form}
    $\mathrm{SU}(2)$群元的一般形式为
    $$\left[\begin{array}{cc}
    a&b\\
    -b^*&a^*
\end{array}\right]\quad |a|^2+|b|^2=1$$
\end{theorem}
\begin{proof}
设
$$A=\left[\begin{array}{cc}
    a&b\\
    c&d
\end{array}\right]$$
$$AA^\dagger=\left[\begin{array}{cc}
    aa^*+bb^*&ac^*+bd^*\\
    a^*c+b^*d&cc^*+dd^*
\end{array}\right]=1\quad \det A=ad-bc=1$$
有
$$|a|^2+|b|^2=1\quad ac^*+bd^*=0\quad ad-bc=1$$
若$a\neq 0$，则
$$c^*=-\frac{bd^*}{a}\quad ad+b\frac{b^* d}{a^*}=1$$
化简后有
$$d=a^*\quad c=-b^*$$
若$a=0$，则$b\neq 0$，同样有
$$d=a^*\quad c=-b^*$$
证毕
\end{proof}
\begin{theorem}
    若$H\in\mathbb{C}^{2\times2}$是二维复零迹厄米矩阵，则$\exp(\i H)\in\mathrm{SU}(2)$。
\end{theorem}
\begin{proof}
    首先有
    $$\exp(\i H)^\dagger \exp(\i H)=\exp(-\i H)\exp(\i H)=I$$
    其次令$A(t)=\exp(\i tH)$，$t\in\mathbb{R}$，再利用公式
    \begin{equation}\label{ddet}
        \frac{\d}{\d t}\det A(t)=\frac{1}{\det A}\mathrm{Tr}\left[A^{-1}\frac{\d A}{\d t}\right]
    \end{equation}
    $$=\frac{1}{\det A}\mathrm{Tr}\left[A(t)^{-1}A(t)\i H\right]=\i\frac{\mathrm{Tr} H}{\det A}=0$$
    因此
    $$\det\exp(\i H)=\det A(1)=\det A(0)=\det I=1$$
    主体证毕。最后我们对于一般的$n$维矩阵，证明式\ref{ddet}。
    $$\frac{\d}{\d t}\det A(t)=\frac{\d}{\d t}\sum_{i_1...i_n}\epsilon_{i_1...i_n}A_{1i_1}...A_{ni_n}=\sum_{s=1}^n\sum_{i_1...i_n}\epsilon_{i_1...i_n}A_{1i_1}...\frac{\d A_{si_s}}{\d t}...A_{ni_n}$$
    又因为
    $$(A^{-1})_{i_s s}=\frac{1}{\det A}\sum_{i_1...i_{s-1}i_{s+1}..i_n}\epsilon_{i_1...i_n}A_{1i_1}..A_{(s-1)i_{(s-1)}}A_{(s+1)i_{(s+1)}}...A_{ni_n}$$
    所以
    $$\frac{\d}{\d t}\det A(t)=\det A\sum_{ij}(A^{-1})_{ij}\frac{\d A_{ji}}{\d t}=\frac{1}{\det A}\mathrm{Tr}\left[A^{-1}\frac{\d A}{\d t}\right]$$
    证毕。
\end{proof}
我们选取二维复零迹厄米矩阵的一组基底，称为Pauli矩阵。
$$\sigma_1=\left[\begin{array}{cc}
    0&1\\
    1&0
\end{array}\right]\quad \sigma_2=\left[\begin{array}{cc}
    0&-\i\\
    \i&0
\end{array}\right]\quad \sigma_3=\left[\begin{array}{cc}
    1&0\\
    0&-1
\end{array}\right]$$
由前面的定理知，任意$\alpha\in\mathbb{R}^3$，$\exp(\i\vec{\alpha}\cdot\vec{\sigma})\in\mathrm{SU}(2)$，下面我们来显式计算它。
$$\exp(\i\vec{\alpha}\cdot\vec{\sigma})=\exp\left(\i\left[\begin{array}{cc}
    \alpha_3&\alpha_1-\i\alpha_2\\
    \alpha+\i\alpha_2&-\alpha_3
\end{array}\right]\right)$$
做对角化处理
$$\left[\begin{array}{cc}
    \alpha_3&\alpha_1-\i\alpha_2\\
    \alpha_1+\i\alpha_2&-\alpha_3
\end{array}\right]=U\left[\begin{array}{cc}
    \alpha&0\\
    0&-\alpha
\end{array}\right]U^\dagger$$
$$U=\left[\begin{array}{cc}
    \dfrac{-\alpha_1+\i\alpha_2}{\sqrt{2\alpha^2-2\alpha\alpha_3}}&\dfrac{-\alpha_1+\i\alpha_2}{\sqrt{2\alpha^2+2\alpha\alpha_3}}\\
    \dfrac{\alpha-\alpha_3}{\sqrt{2\alpha^2-2\alpha\alpha_3}}&\dfrac{-\alpha-\alpha_3}{\sqrt{2\alpha^2+2\alpha\alpha_3}}
\end{array}\right]$$
$$\exp(\i\vec{\alpha}\cdot\vec{\sigma})=\left[\begin{array}{cc}
    \cos\alpha+\i\dfrac{\alpha_3}{\alpha}\sin\alpha&\dfrac{\alpha_1-\i\alpha_2}{\alpha}\i\sin\alpha\\
    \dfrac{\alpha_1+\i\alpha_2}{\alpha}\i\sin\alpha&\cos\alpha-\i\dfrac{\alpha_3}{\alpha}\sin\alpha
\end{array}\right]\quad\alpha=\sqrt{\alpha_1^2+\alpha_2^2+\alpha_3^2}$$

由定理\ref{SU2_form}知，$\exp(\i\vec{\alpha}\cdot\vec{\sigma})$包含了所有SU(2)的元素。$\vec{\sigma}$称作SU(2)的生成元。
\subsection{角动量耦合}
CG系数
\begin{align}\label{CG-defination}
  |jm\rangle = \sum_{m_1,m_2}C_{j_1 j_2}(jm;m_1m_2)|j_1 m_1;j_2 m_2\rangle  
\end{align}

\begin{theorem}[正交关系]
\begin{align*}
    \sum_{m_1,m_2} C^*_{j_1 j_2}(j'm';m_1m_2)C_{j_1 j_2}(jm;m_1m_2)= \delta_{j'j}\delta_{m'm}
\end{align*}
\begin{align*}
    \sum_{j,m}C^*_{j_1,j_2}(jm;m_1'm_2')C_{j_1,j_2}(jm;m_1m_2)= \delta_{m_1'm_1}\delta_{m_2'm_2}
\end{align*}
\end{theorem}
\begin{proof}
    利用两套基矢的正交归一性，按CG系数定义\ref{CG-defination}，
    \begin{align*}
        &\langle j'm'| jm\rangle \\
        &= \sum_{\substack{m_1,m_2\\m_1',m_2'}} C^*_{j_1 j_2}(j'm';m_1'm_2')C_{j_1 j_2}(jm;m_1m_2)\langle j_1 m_1';j_2 m_2'|j_1 m_1;j_2 m_2\rangle\\
&= \sum_{\substack{m_1,m_2\\m_1',m_2'}} C^*_{j_1 j_2}(j'm';m_1'm_2')C_{j_1 j_2}(jm;m_1m_2)\delta_{m_1' m_1}\delta_{m_2' m_2}\\
&= \sum_{m_1,m_2} C^*_{j_1 j_2}(j'm';m_1m_2)C_{j_1 j_2}(jm;m_1m_2)\\
&= \delta_{j'j}\delta_{m'm}
    \end{align*}
    反过来展开：
    \begin{align*}
        |j_1 m_1;j_2 m_2\rangle = \sum_{j,m}\langle jm|j_1 m_1;j_2 m_2 \rangle|jm\rangle= \sum_{j,m}C^*_{j_1,j_2}(jm;m_1m_2)|jm\rangle
    \end{align*}
    同样操作
    \begin{align*}
        &\langle j_1 m_1';j_2 m_2'|j_1 m_1;j_2 m_2\rangle\\
        &= \sum_{\substack{j,m\\j',m'}}C^*_{j_1,j_2}(jm;m_1m_2)C_{j_1,j_2}(j'm';m_1'm_2')\langle j'm'|jm\rangle\\
&= \sum_{\substack{j,m\\j',m'}}C^*_{j_1,j_2}(jm;m_1m_2)C_{j_1,j_2}(j'm';m_1'm_2')\delta_{j'j}\delta_{m'm}\\
&= \sum_{j,m}C^*_{j_1,j_2}(jm;m_1m_2)C_{j_1,j_2}(jm;m_1'm_2')\\
&= \delta_{m_1'm_1}\delta_{m_2'm_2}
    \end{align*}
\end{proof}
\begin{theorem}
Racah表示式
\begin{align*}
    &C_{j_1j_2}(j m;m_1 m_2) \\
    &= \delta_{m1+m2,m}\sqrt{(2j+1)\frac{(j_1+j_2-j)!(j_2+j-j_1)!(j+j_1-j_2)!}{(j_1+j_2+j+1)!}}\\
    &\times\sqrt{(j_1+m_1)!(j_1-m_1)!(j_2+m_2)!(j_2-m_2)!(j+m)!(j-m)!}\\
    &\times\sum_{k}\frac{(-1)^{k}}{k!(j_1+j_2-j-k)!(j_1-m_1-k)!(j_2+m_2-k)!(j-j_1-m_2+k)!(j-j_2+m_1+k)!}
\end{align*}
\end{theorem}
\begin{proof}
用升降算符作用：
\begin{align*}
    L_{\pm}|jm\rangle  = \sqrt{(j\mp m)(j\pm m +1)}\hbar|j,m\pm1\rangle
\end{align*}
\begin{align*}
    &\sum C_{j_1j_2}(j m;m_1 m_2)(L_{1\pm} + L_{2\pm})|j_1 m_1;j_2 m_2\rangle\\
    &=\sum C_{j_1j_2}(j m;m_1 m_2)\left(\begin{matrix}\sqrt{(j_1\mp m_1)(j_1\pm m_1 +1)}\hbar|j_1, m_1\pm1;j_2 m_2\rangle\\
    +\sqrt{(j_2\mp m_2)(j_2\pm m_2 +1)}\hbar|j_1 m_1;j_2,m_2\pm1\rangle\end{matrix}\right)
\end{align*}
比较系数
\begin{align*}
    &\sqrt{(j\mp m)(j\pm m +1)}C_{j_1j_2}(j, m\pm1;m_1 m_2)\\
    &= \sqrt{(j_1\pm m_1)(j_1\mp m_1 +1)}C_{j_1j_2}(jm;m_1\mp1, m_2)\\ 
    &+ \sqrt{(j_2\pm m_2)(j_2\mp m_2 +1)}C_{j_1j_2}(j m;m_1, m_2\mp1)
\end{align*}
先看最高的$m$使得$L_{+}|jm\rangle = 0$，即$m=j$：
\begin{align*}
    &\sum_{m_1 m_2}C_{j_1j_2}(j j;m_1 m_2) \left(\begin{matrix}
    \sqrt{(j_1-m_1)(j_1+m_1+1)}\hbar|j_1, m_1+1;j_2 m_2\rangle\\
    + \sqrt{(j_2-m_2)(j_2+m_2+1)}|j_1 m_1;j_2,m_2+1\rangle\end{matrix}\right) = 0
\end{align*}
\begin{align*}
    &C_{j_1j_2}(j j;m_1-1,j-m_1)\sqrt{(j_1+m_1)(j_1-m_1+1)} \\
    &+C_{j_1j_2}(j j;m_1,j-m_1-1)\sqrt{(j_2+m_2)(j_2-m_2+1)} = 0
\end{align*}
设$m_2 = j_2 - k, m_1 = j - j_2 + k,C_k = C_{j_1j_2}(jj;j-j_2+k,j_2-k)$，
\begin{align*}
    \frac{C_k}{C_{k-1}} = -\sqrt{\frac{(j_1+m_1)(j_1-m_1+1)}{(j_2-m_2)(j_2+m_2+1)}}=-\sqrt{\frac{(j_1+j - j_2 + k)(j_1-j + j_2 - k+1)}{k(2j_2-k+1)}}
\end{align*}
\begin{align*}
    C_k = C_0 (-1)^k \sqrt{\frac{(2j_2-k)!(j_1+j-j_2+k)!(j_1-j+j_2)!}{k!(2j_2)!(j_1-j+j_2-k)!}}
\end{align*}
对$|jj\rangle$归一化，
\begin{align*}
    &|C_0|^2\sum_{k=0}^{j_1+j_2-j}\frac{(2j_2-k)!(j_1-j_2+j+k)!}{k!(j_1+j_2-j-k)!} \\
&= (j_2-j_1+j)!(j_1-j_2 + j)!\sum_{k=0}^{j_1+j_2-j}\left(\begin{matrix}2j_2-k\\j_2-j_1+j\end{matrix}\right)\left(\begin{matrix}j_1-j_2+j+k\\j_1-j_2+j\end{matrix}\right)\\
&=1
\end{align*}
使用母函数计算求和：
\begin{align*}
    F(x) = \sum_{k}\left(\begin{matrix}j_1-j_2+j+k\\j_1-j_2+j\end{matrix}\right) x^k = \frac{1}{(1-x)^{j_1-j_2+j+1}}
\end{align*}
\begin{align*}
    G(x) = \sum_{k}\left(\begin{matrix}j_2+j-j_1+k\\j_2-j_1+j\end{matrix}\right)x^k =\frac{1}{(1-x)^{j_2-j_1+j+1}}
\end{align*}
待求和项恰好为两者卷积：
\begin{align*}
    &\sum_{k=0}^{j_1+j_2-j}\left(\begin{matrix}2j_2-k\\j_2-j_1+j\end{matrix}\right)\left(\begin{matrix}j_1-j_2+j+k\\j_1-j_2+j\end{matrix}\right) = \sum_{k=0}^{j_1+j_2-j} f_k g_{j_1+j_2-j-k}\\
    &= [x^{j_1+j_2-j}]F(x)G(x) = \left(\begin{matrix}j+j_1+j_2+1\\j_1+j_2-j\end{matrix}\right)
\end{align*}
其中$[x^k]f(x)$记号表示$f(x)$中$x^k$项的系数.\\
 Condon–Shortley 约定要求
 $$
 C_{j_1j_2}(jj;j_1,j-j_1) > 0
 $$
 即在$C_k$前加$(-1)^{j-j_1-j_2}$.
 于是
\begin{align*}
    &C_{j_1j_2}(j j;m_1,m_2) \\
    &= \delta_{m_1+m_2,j}(-1)^{j-j_1-m_2}\sqrt{\frac{(j_1+j_2-j)!(2j+1)!}{(j_2-j_1+j)!(j_1-j_2 + j)!(j+j_1+j_2+1)!}}\sqrt{\frac{(j_1+m_1)!(j_2+m_2)!}{(j_1-m_1)!(j_2-m_2)!}}
\end{align*}

接下来从$|jj\rangle$下降到一般的$|jm\rangle$态，将$L_{-}$重复作用：
\begin{align*}
    &|jm\rangle = \sqrt{\frac{(j+m)!}{(2j)!(j-m)!}}J_{-}^{j-m}|jj\rangle\\
&= \sqrt{\frac{(j+m)!}{(2j)!(j-m)!}}\sum\left(\begin{matrix}j-m\\ s\end{matrix}\right)J_{1-}^{s} J_{2-}^{j-m-s}|jj\rangle\\
&= \sqrt{\frac{(j+m)!}{(2j)!(j-m)!}}\sum_s\left(\begin{matrix}j-m\\ s\end{matrix}\right)J_{1-}^{s} J_{2-}^{j-m-s}\sum_{m_1,m_2} C_{j_1j_2}(j j;m_1',m_2') |j_1,m_1';j_2,m_2'\rangle\\
&= \sqrt{\frac{(j+m)!}{(2j)!(j-m)!}}\sum_s\left(\begin{matrix}j-m\\ s\end{matrix}\right)\sqrt{\frac{(j_1 + m_1')!(j_1 -m_1' +s)!}{(j_1 - m_1')!(j_1 +m_1'-s)!}\frac{(j_2 + m_2')!(j_2+j-m_2'-m-s)!}{(j_2-m_2')!(j_2 - j +m_2'+m+s)!}} \\
&\times\sum_{m_1',m_2'} C_{j_1j_2}(j j;m_1',m_2')|j_1,m_1'-s;j_2,m_2'-j+m+s\rangle
\end{align*}
针对希望计算的$m_1 m_2$，有
\begin{align*}
    m_1' = m_1+s, m_2' = m_2+j-m-s
\end{align*}
\begin{align*}
    &C_{j_1 j_2}(jm;m_1m_2)\\
    &= \delta_{m_1+m_2,m}\sum_{s} \sqrt{\frac{(j+m)!}{(2j)!(j-m)!}}\frac{(j-m)!}{s!(j - m - s)!}\\
    &\times \sqrt{\frac{(j_1 + m_1 + s)!(j_1 -m_1)!}{(j_1 - m_1-s)!(j_1 +m_1)!}\frac{(j_2 + m_2+j-m-s)!(j_2-m_2)!}{(j_2-m_2-j+m+s)!(j_2+m_2)!}}\\
    &\times (-1)^{j_1 - m_2+m+s}\sqrt{\frac{(j_1+j_2-j)!(2j+1)!}{(j_2-j_1+j)!(j_1-j_2 + j)!(j+j_1+j_2+1)!}}\sqrt{\frac{(j_1+m_1+s)!(j_2+m_2+j-m-s)!}{(j_1-m_1-s)!(j_2-m_2-j+m+s)!}}\\
    &= \delta_{m_1+m_2,m} \sqrt{(2j+1)}\sqrt{(j+m)!(j-m)!(j_1+m_1)!(j_1-m_1)!(j_2+m_2)!(j_2-m_2)!}\\
    &\times \sqrt{\frac{(j_1+j_2-j)!}{(j_2-j_1+j)!(j_1-j_2 + j)!(j+j_1+j_2+1)!}}\\
    &\times \sum_{s} \frac{(-1)^{-j_1 - m_2+m+s} (j_1 + m_1 + s)!(j_2+j-m_1-s)!}{s!(j - m - s)!(j_1 - m_1-s)!(j_2-j+m_1+s)!(j_1 + m_1)!(j_2 + m_2)!}
\end{align*}
这被称为Edmonds形式，可惜此结果不对称，Racah在他的论文中将其变换为了熟知的结果%\footcite{PhysRev.62.438}
.也可参见喀兴林《高等量子力学》5.23.我们接下来按照喀兴林中习题的引导证明.
\end{proof}
\begin{lemma}
证明
\begin{align}\label{lemma1}
    \sum_{r} \binom{n}{r}\binom{m}{t-r}=\binom{n+m}{t}
\end{align}
对任意整数 (n,m) 都成立.求和要求所有阶乘非负，下同.
\end{lemma}
\begin{proof}
    我们已在前使用母函数证明了这一点.
\end{proof}
\begin{lemma}
证明
    \begin{align}\label{lemma2_1}
        \frac{a!}{b!c!}=\sum_{r=0}^\infty\frac{(a-b)!(a-c)!}{(a-b-r)!(a-c-r)!(b+c-a+r)!r!}
    \end{align}
以及
\begin{align}\label{lemma2_2}
   \frac{(n+m)!}{n!m!s!(n+m-s)!}=\sum_{r=0}^\infty\frac{1}{(n-r)!(s-r)!(m-s+r)!r!}
\end{align}
\end{lemma}
\begin{proof}
在\ref{lemma1}中代入$n+m=a,m=b,n+m-t=c$:
\begin{align*}
    \sum_{r=0}^\infty \binom{a-b}{r}\binom{b}{a-c-r}=\binom{a}{a-c}
\end{align*}
二项系数写为阶乘：
\begin{align*}
    \sum_{r=0}^\infty \frac{(a-b)!}{r!(a-b-r)!} \frac{b!}{(a-c-r)!(b+c-a+r)!}=\frac{a!}{(a-c)!c!}
\end{align*}
两边同乘 $\dfrac{(a-c)!}{b!}$，证得
\begin{align*}
     \frac{a!}{b!c!}=\sum_{r=0}^\infty\frac{(a-b)!(a-c)!}{(a-b-r)!(a-c-r)!(b+c-a+r)!r!}
\end{align*}
再证第二式，令$a = n + m, b=m,a-c=s$，代入：
\begin{align*}
    \frac{(n+m)!}{m!(n+m-s)!}=\sum_{r=0}^\infty
\frac{n!s!}{(n-r)!(s-r)!(m-s+r)!r!}
\end{align*}
两边同除$n!m!s!$，证得
\begin{align*}
    \frac{(n+m)!}{n!m!s!(n+m-s)!}=\sum_{r=0}^\infty\frac{1}{(n-r)!(s-r)!(m-s+r)!r!}
\end{align*}
\end{proof}
\begin{lemma}
    \begin{align}\label{lemma3_1}
        \frac{(s-n)!(s-t)!}{n!t!(s-n-t)!}=\sum_{r=0}^\infty(-1)^r\frac{(s-r)!}{r!(n-r)!(t-r)!}
    \end{align}
    以及
    \begin{align}\label{lemma3_2}
        \frac{(a+b+1)!(a-c)!(b-d)!}{(c+d)!(a+b-c-d+1)!}=\sum_{s=0}^\infty\frac{(a+s)!(b-s)!}{(c+s)!(d-s)!}
    \end{align}
\end{lemma}
\begin{proof}
    \begin{align*}
        \sum_{r=0}^\infty \binom{n}{r}\binom{-(s-t+1)}{t-r}=\binom{n-(s-t+1)}{t}
    \end{align*}
    \begin{align*}
        \binom{-(s-t+1)}{t-r}=(-1)^{t-r}\frac{(s-r)!}{(t-r)!(s-t)!}
    \end{align*}
    因此左边为
\[
\sum_{r=0}^\infty\frac{n!}{r!(n-r)!}(-1)^{t-r}\frac{(s-r)!}{(t-r)!(s-t)!}
\]
在右边再用一次负上指标公式：

\[\binom{-(s-n-t+1)}{t}=(-1)^t\frac{(s-n)!}{t!(s-n-t)!}.
\]
\[ \frac{(-1)^t n!}{(s-t)!} \sum_{r=0}^\infty (-1)^r \frac{(s-r)!}{r!(n-r)!(t-r)!}=(-1)^t\frac{(s-n)!}{t!(s-n-t)!}
\]
移项即得
\[
\frac{(s-n)!(s-t)!}{n!t!(s-n-t)!}=\sum_{r=0}^\infty(-1)^r\frac{(s-r)!}{r!(n-r)!(t-r)!}
\]
接下来证明第二式，
\begin{align*}
    &\frac{(a+s)!(b-s)!}{(c+s)!(d-s)!}\\
    &=(a-c)!(b-d)!\binom{a+s}{c+s}\binom{b-s}{d-s}\\
    &=(a-c)!(b-d)!(-1)^{c+d}\binom{c-a-1}{c+s}\binom{d-b-1}{d-s}
\end{align*}
使用\ref{lemma1}，
\begin{align*}
    \sum\binom{c-a-1}{c+s}\binom{d-b-1}{d-s} = \binom{c+d-a-b-2}{c+d} = (-1)^{c+d}\binom{a+b+1}{c+d}
\end{align*}
于是
\begin{align*}
    &\sum\frac{(a+s)!(b-s)!}{(c+s)!(d-s)!} \\
    &= (a-c)!(b-d)!\binom{a+b+1}{c+d}\\
    &= \frac{(a+b+1)!(a-c)!(b-d)!}{(c+d)!(a+b-c-d+1)!}
\end{align*}
\end{proof}
\begin{theorem}
    Edmonds形式与Racah形式等价.
\end{theorem}
\begin{proof}
在\ref{lemma2_2}中代入$(n,m,s)=(j-j_1+j_2,j_1-m_1-s,j_2+m_2)$得到：
\begin{align*}
    &\frac{(j+j_2-m_1-s)!} {(j_1-m_1-s)!(j-m-s)!(j-j_1+j_2)!(j_2+m-m_1)!}\\
    &=\sum_u\frac{1}{u!(j_1-j_2-m-s+u)!(j-j_1+j_2-u)!(j_2+m_2-u)!}
\end{align*}
在\ref{lemma3_1}中代入$(n,t,s,r)=(j_1-j_2-m+u,j_1-j-m_2+u,2j_1-j_2-m_2-u,j_1-j_2-m-u-s)$得到：
\begin{align*}
    &\sum_s (-1)^s \frac{(j_1+m_1+s)!} {s!(j_2-j+m_1+s)!(j_1-j_2-m+u-s)!}\\
&=(-1)^{j_1-j_2-m+u}\frac{(j_1+m_1)!(j+j_1-j_2)!}{(j_1-j_2-m+u)!(j_1-j-m_2+u)!(j+m-u)!}
\end{align*}
\begin{align*}
    &\sum_{s} \frac{(-1)^{-j_1 - m_2+m+s} (j_1 + m_1 + s)!(j_2+j-m_1-s)!}{s!(j - m - s)!(j_1 - m_1-s)!(j_2-j+m_1+s)!(j_2-j_1+j)!(j_1-j_2 + j)!(j_1 + m_1)!(j_2 + m_2)!}\\
    &=\sum_{s}\frac{(j_2+j-m_1-s)!}{(j - m - s)!(j_1 - m_1-s)!(j_2-j_1+j)!(j_2 + m_2)!}\\
    &\times(-1)^{s}\frac{(j_1 + m_1 + s)!}{s!(j_2-j+m_1+s)!}\cdot \frac{1}{(j_1-j_2 + j)!(j_1 + m_1)!}\\
    &=\sum_{s}\sum_{u} \frac{1}{u!(j-j_1+j_2-u)!(j_2+m_2-u)!}\cdot \frac{1}{(j_1-j_2 + j)!(j_1 + m_1)!}\\
    &\times(-1)^{s}\frac{(j_1 + m_1 + s)!}{s!(j_2-j+m_1+s)!(j_1-j_2-m-s+u)!}\\
    &=\sum_{u}(-1)^{-j_1+m_1}\frac{(-1)^{j_1-j_2-m+u}}{u!(j-j_1+j_2-u)!(j_2+m_2-u)!(j_1-j_2-m+u)!(j_1-j-m_2+u)!(j+m-u)!}
\end{align*}
做换元$k=j_1-j-m_2+u$：
\begin{align*}
    =\sum_{k}\frac{(-1)^{k}}{k!(j_1+j_2-j-k)!(j_1-m_1-k)!(j_2+m_2-k)!(j-j_2+m_1+k)!(j-j_1-m_2+k)!}
\end{align*}
\end{proof}
好消息是实践中$j_1,j_2$总是取较小的几个值，我们不需要计算这个复杂的式子，不过计算步骤和上面完全相同：先计算最高的$m=j$，然后用降算符得到一般的$|lm\rangle$.

\section{有心力场}
\subsection{代数法解氢原子}
哈密顿量
$$
H = \frac{p^2}{2\mu} - \frac{\kappa}{r}
$$
隆格-楞茨矢量
$$
\vec R = \frac{1}{2\mu\kappa}\left(\vec p\times \vec L - \vec L \times \vec p\right) - \frac{\vec r}{r}
$$
利用
\begin{align*}
&\vec p\times \vec L + \vec L \times \vec p \\
&= \epsilon_{ijk}p_i \epsilon_{lmj}r_l p_m + \epsilon_{ijk}\epsilon_{lmi}r_l p_m p_j\\
&= -p_i r_i p_k + p_i r_k p_i + r_j p_k p_j - r_k p_j p_j\\
&= (-p_i r_i+r_i p_i) p_k + (p_i r_k - r_k p_i)p_i = 3i\hbar p_k - i \hbar \delta_{ik}p_i = 2i\hbar p_k
\end{align*}
化为
\begin{align*}
\vec R = \frac{1}{\mu \kappa}(\vec p \times \vec L - i\hbar \vec p) - \frac{\vec r}{r}
\end{align*}
\begin{theorem}
$[R_i, H] = 0$
\end{theorem}
\begin{proof}
\begin{align*}
[R,\frac{1}{r}] &=  \frac{1}{\mu \kappa}([\vec p, \frac{1}{r}] \times \vec L + \vec p \times [\vec L, \frac{1}{r}] - i\hbar [\vec p, \frac{1}{r}])\\
&=\frac{1}{\mu\kappa}\left(i\hbar \frac{\vec r}{r^3}\times \vec L + \hbar^2 \frac{\vec r}{r^3}\right)\\
&= \frac{1}{\mu\kappa}\left(i\hbar \frac{1}{r^3} \vec r(\vec r\cdot \vec p) - \frac{1}{r}\vec p + \hbar^2 \frac{\vec r}{r^3}\right)
\end{align*}
\begin{align*}
[R, p^2] &= -[\frac{\vec r}{r}, p^2]\\
&= - [\frac{\vec r}{r}, \vec p] \cdot \vec p - \vec p \cdot[\frac{\vec r}{r}, \vec p]\\
&= 2 i\hbar\left(\frac{\vec r \vec r}{r^3} \cdot \vec p - \frac{1}{r}\vec p + \vec p \cdot \frac{\vec r \vec r}{r^3}  - \vec p \frac{1}{r}\right)\\
&= i\hbar\left(\frac{\vec r \vec r}{r^3} \cdot \vec p - \frac{1}{r}\vec p\right) + 2\hbar^2 \frac{\vec r}{r^3}
\end{align*}
\begin{align*}
[R,H] = \frac{1}{2\mu}[R,p^2] -\kappa [R, \frac{1}{r}] = 0
\end{align*}
\end{proof}
\begin{theorem}
$[R_i, R_j] = 0$
\end{theorem}
\begin{proof}
\begin{align*}
[R_i, R_j] &= \frac{1}{\mu^2\kappa^2}[(\vec p \times \vec L)_i,(\vec p \times \vec L)_j] \\
&- \frac{i\hbar}{\mu^2\kappa^2} ([(\vec p \times \vec L)_i, p_j] - (i \leftrightarrow j))\\
&+ \frac{i\hbar}{\mu\kappa} ([p_i, \frac{r_j}{r}] - (i \leftrightarrow j))\\
&- \frac{1}{\mu\kappa} ([(\vec p \times \vec L)_i, \frac{r_j}{r}] - (i \leftrightarrow j))
\end{align*}
其中$(i \leftrightarrow j)$表示将$ij$指标交换，这提示对于后三项仅需考虑反对称部分.依次计算：
\begin{align*}
   & [(\vec p \times \vec L)_i, p_j] - (i \leftrightarrow j)\\
&= \epsilon_{ikl} p_k [L_l, p_j] - (i \leftrightarrow j)\\
&= i\hbar\epsilon_{ikl}\epsilon_{ljm} p_k p_m - (i \leftrightarrow j)\\
&= i\hbar(\delta_{ij} p_k p_k - p_i p_j) - (i \leftrightarrow j)\\
&= 0
\end{align*}
\begin{align*}
    [p_i, \frac{r_j}{r}] -(i \leftrightarrow j) = i\hbar\left(\frac{r_i r_j}{r^3} - \frac{\delta_{ij}}{r}\right) -(i \leftrightarrow j) = 0
\end{align*}
\begin{align*}
   & [(\vec p \times \vec L)_i, \frac{r_j}{r}] - (i \leftrightarrow j)\\
&= \epsilon_{ikl} \left([p_k, \frac{r_j}{r}]L_l + p_k[L_l, \frac{r_j}{r}]\right) - (i \leftrightarrow j)\\
&=\epsilon_{ikl} \left(i\hbar \left(\frac{r_j r_k}{r^3} - \frac{\delta_{jk}}{r}\right)L_l + p_k i\hbar \epsilon_{ljm}\frac{r_m}{r}\right) - (i \leftrightarrow j) \\
&= i\hbar\left(\epsilon_{ikl}\frac{1}{r^3}r_k r_j L_l - \epsilon_{ijl}\frac{1}{r}L_l + \delta_{ij}p_k \frac{r_k}{r} - p_j \frac{r_i}{r}\right) - (i \leftrightarrow j) \\
&= i\hbar \frac{r_j r_k}{r^3}\epsilon_{ikl}\epsilon_{mnl}r_m p_n - i\hbar p_j\frac{r_i}{r} - (i \leftrightarrow j) - 2i\hbar \epsilon_{ijl}\frac{1}{r}L_l\\
&= \frac{i\hbar}{r^3} (r_i r_j r_k p_k - r^2 r_j p_i)- i\hbar p_j\frac{r_i}{r} - (i \leftrightarrow j) - 2i\hbar \epsilon_{ijl}\frac{1}{r}L_l\\
&= - 2i\hbar \epsilon_{ijl}\frac{1}{r}L_l
\end{align*}
\begin{align*}
    &[(\vec p \times \vec L)_i,(\vec p \times \vec L)_j] \\
&= \epsilon_{ikl}\epsilon_{jmn}(p_k [L_l, p_m] L_n + p_m [L_l, L_n] + p_m[p_k, L_n]L_l)\\
&= i\hbar \epsilon_{ikl}\epsilon_{jmn} (p_k \epsilon_{lms}p_s L_n + p_k p_m \epsilon_{lns} L_s - p_m \epsilon_{nks} p_s L_l)\\
\end{align*}
拆成三项计算：
\begin{align*}
&\epsilon_{ikl}\epsilon_{jmn}\epsilon_{lms}p_k p_s L_n\\
&=\epsilon_{ikl}(\delta_{jl}\delta_{ns} - \delta_{js}\delta_{nl})p_k p_s L_n \\
&= -\epsilon_{ikj} p_k p_n L_n - \epsilon_{ikl} p_k p_j L_l\\
&= -\epsilon_{ikl}p_k p_j L_l
\end{align*}
其中利用$\vec p \cdot \vec L = 0$.
\begin{align*}
&\epsilon_{ikl}\epsilon_{jmn}\epsilon_{lns}p_k p_m L_s\\
&\epsilon_{ikl}(\delta_{jl}\delta_{ms}-\delta_{js}\delta_{ml})p_k p_m L_s\\
&= -\epsilon_{ikj} p_k p_m L_m + \epsilon_{ikl}p_k p_l L_j\\
&= \epsilon_{ikl}p_k p_l L_j = 0
\end{align*}
其中利用$\epsilon_{ikl}$反对称，而$p_k p_l$对称，因此求和为$0$.
\begin{align*}
&\epsilon_{ikl}\epsilon_{jmn}\epsilon_{nks}p_m p_s L_l\\
&=-\epsilon_{ikl}(\delta_{jk}\delta_{ms} - \delta_{js}\delta_{mk})p_m p_s L_l\\
&= -\epsilon_{ijl} p_m p_m L_l + \epsilon_{ikl} p_k p_j L_l\\
&= -\epsilon_{ijl} p^2 L_l + \epsilon_{ikl} p_k p_j L_l
\end{align*}
总和
\begin{align*}
    [(\vec p \times \vec L)_i,(\vec p \times \vec L)_j] = -i\hbar \epsilon_{ijk}p^2 L_k
\end{align*}
证得
\begin{align*}
    [R_i, R_j] = - i\hbar \left(\frac{p^2}{\mu^2\kappa^2} - 2i\hbar\frac{1}{\mu\kappa}\frac{1}{r}\right) \epsilon_{ijk}L_k = - \frac{2i\hbar}{\mu\kappa^2}H\epsilon_{ijk}L_k
\end{align*}
\end{proof}
\begin{theorem}
$R^2 = \frac{2H}{\mu\kappa^2}(L^2 + \hbar^2) + 1$
\end{theorem}
\begin{proof}
\begin{align*}
    &\epsilon_{ijk}p_j L_k \epsilon_{ilm} p_l L_m\\
    &= (\delta_{jl}\delta_{km} - \delta_{jm}\delta_{kl})p_j L_k p_l L_m \\
    &= p_j L_k p_j L_k - p_j L_k p_k L_j \\
    &= p_j (p_j L_k + i\hbar \epsilon_{kjl}p_l)L_k + 0\\
    &= p^2 L^2 + i\hbar \epsilon_{kjl}p_j p_l L_k \\
    &=p^2L^2
\end{align*}
\begin{align*}
    &\epsilon_{ijk}p_j L_k p_i= -\epsilon_{ijk}L_j p_k p_i + 2i\hbar p_i p_i= 2i\hbar p^2
\end{align*}
\begin{align*}
    p_i \epsilon_{ijk}p_jL_k = 0
\end{align*}
\begin{align*}
    \epsilon_{ijk}p_j L_k r_i = (-\epsilon_{ijk} L_j p_k + 2i\hbar p_i)r_i = -L_j \epsilon_{kij} p_k r_i + 2i\hbar p_i r_i = - L^2 + 2i\hbar p_i r_i
\end{align*}
\begin{align*}
    r_i\epsilon_{ijk}p_j L_k = L_k L_k
\end{align*}
证得
\begin{align*}
    R^2 &= \frac{1}{\mu^2\kappa^2} (p^2 L^2 - \hbar^2 p^2 - i\hbar 2i\hbar p^2) - \frac{1}{\mu\kappa}\left((L^2+2i\hbar p_i r_i)\frac{1}{r} - i\hbar p_i r_i \frac{1}{r} -\frac{1}{r}(L^2 - i\hbar r_i p_i)\right) +1\\
    &= \frac{1}{\mu^2\kappa^2} (p^2 L^2 + \hbar^2 p^2) - \frac{1}{\mu\kappa} \left(\frac{2}{r}L^2 + i\hbar p_i \frac{r_i}{r} - i\hbar \frac{1}{r}r_i p_i\right) + 1\\
    &= \frac{p^2}{\mu^2\kappa^2}(L^2 + \hbar^2) - \frac{2}{\mu\kappa}\frac{1}{r}(L^2+\hbar^2) + 1\\
    &=\frac{2}{\mu\kappa^2}H (L^2+\hbar^2) +1
\end{align*}
\end{proof}
\begin{theorem}
    $R\cdot L = L\cdot R = 0$
\end{theorem}
\begin{proof}
    \begin{align*}
        R_i L_i = \frac{1}{\mu\kappa} (\epsilon_{ijk}p_jL_k L_i - i\hbar p_i L_i) - \frac{1}{r}r_i \epsilon_{ijk}r_j p_k = 0
    \end{align*}
    \begin{align*}
        [L_i, R_i] = i\hbar \epsilon_{ijk} \delta_{ij} R_k = 0 \rightarrow R\cdot L = L\cdot R = 0
    \end{align*}
\end{proof}
在束缚态能级$E$对应的子空间中，定义
$$
A = \sqrt{-\frac{\mu\kappa^2}{2E}}R
$$
有
$$
\begin{cases}
[L_i, L_j] = i\hbar \epsilon_{ijk}L_k,\\
[L_i, A_j] = i\hbar \epsilon_{ijk}A_k,\\
[A_i, A_j] = i\hbar \epsilon_{ijk}L_k
\end{cases}
$$
这构成了SO4对称性.

定义
$$
A_{\pm} = \frac{1}{2}\left( L \pm A \right)
$$
有
$$
[A_{\pm i}, A_{\pm j}]\\ 
= \frac{1}{4} \left([L_i, L_j] + [A_i, A_j] \pm [L_i, A_j] \pm [A_i, L_j] \right)\\
= i\epsilon_{ijk} A_{\pm k} 
$$
这说明$A_{\pm}$各自构成一套SO3群.类似于角动量，$A_{\pm}^2$的本征值为$a_{\pm}(a_{\pm}+1)\hbar^2$，由于$A_{+}^2 = A_{-}^2$，$a_+ = a_- = a$,

$$
A^2 = (A_{+} - A_{-})^2 = - (A_{+} + A_{-})^2 - \hbar^2 - \frac{\mu\kappa^2}{2E}
$$
$$
(4 a(a+1) + 1) \hbar^2 = -\frac{\mu\kappa^2}{2E}\rightarrow E = -\frac{\mu\kappa^2}{2(2a+1)^2}
$$
$n = 2a+1,\kappa = \frac{Ze^2}{4\pi\epsilon_0}$,证得
$$
E = - \frac{\mu Z^2 e^4}{32\pi^2\epsilon_0^2n^2}
$$

\end{document}